\subsection{Az "M" nyelv (Matlab) jellegzetességei: változók, vektorok és mátrixok, feltételes végrehajtás, ciklusok, számtani sorozatok, függvény definíció, diagram rajzolás.}

\textbf{Az ,,M'' nyelv jellegzetességei:}
\begin{itemize}
    \item interpreter ,,.m'' fájlt tud lefuttatni, függvényként indítja el
    \item \verb|;| -vel zárva a sort nem jeleníti meg az eredményt
    \item változó kiírása: \verb|disp();|
    \item komment jele: \verb|%|
    \item \verb|%%|: szekció -- lehet csak szekciót futtatni
    \item változókat nem kell deklarálni, értékadásnál kapnak típust
    \item numerikus értékek double pontosak, komplex számok is vannak: \verb|i| és \verb|j| is a képzetes rész
    \item string konstans jele az egyes és kettes macskaköröm is
    \item kettes macskakörömben lévő stringek összefűzhetők \verb|+| jellel
    \item \verb|num2str()|: szám -> szöveg konverzió
    \item \textbf{változó terület neve workspace}
    \begin{itemize}
        \item értékeit megtartja futtatások közt
        \item \verb|clear| -> törlés
        \item kimenthető és visszatölthető -- \verb|.mat| fájlba
    \end{itemize}
    \item \textbf{feltételes utasítás}
    \begin{itemize}
        \item feltétel nincs zárójelben
        \item végére \verb|end| kell
        \item 0 a hamis
    \end{itemize}
    \item létezik \verb|&|, \verb|&&|: ,,és'', \verb+|+ , \verb+||+: ,,vagy''
    \begin{itemize}
        \item Xor műveleti jel nincs
        \item \verb|xor()|, \verb|bitxor()| függvények vannak
    \end{itemize}
    \item \verb|^| a hatványozás jele
    \item \textbf{függvény definíció}
    \begin{itemize}
        \item vagy külön \verb|.m| fájlban
        \item vagy a programunk végén
    \end{itemize}
    \item nincs hátultesztelt ciklus
    \item \textbf{számtani sorozat készítés}
    \begin{itemize}
        \item for ciklushoz is ezt használjuk
        \item kezdőérték : végérték
        \item kezdőérték : növekmény : végérték
        \item \verb|linspace(kezdőérték, végérték, darab)| függvény
        \begin{itemize}
            \item megadott tartományt megadott darab részre osztja
        \end{itemize}
    \end{itemize}
    \item \textbf{műveletek vektorokkal}
    \begin{itemize}
        \item ha a vektor elemein akarunk műveletet végezni: ponttal jelöljük
        \item pl: \verb|C = B ./ a;|  \% B vektor elemeit egyesével osztottuk le 'a' konstanssal
    \end{itemize}
    \item \textbf{rajzolás 2D}
\begin{minted}[breaklines]{matlab}
figure; % ábra window
clf; % grafika törlése
hold on; % minden plot ugyanoda
x=0.1:0.1:4*pi; % kezd:növekmény:vég
y1=sin(x);
plot(x,y1,'green');
y2=sin(x)./x; % elemenkénti osztás
plot(x,y2,'blue');
\end{minted}
    \item \textbf{függvény automatikus rajzolás}
    \begin{itemize}
        \item \verb|fplot()|
    \end{itemize}
    \item \textbf{kirajzolt ablak skálázása}
    \begin{itemize}
        \item \verb|grid();|
    \end{itemize}
    \item \textbf{címsor}
    \begin{itemize}
        \item \verb|title(string)|
    \end{itemize}
    \item \textbf{méretarányos skálaosztás a tengelyeken}
    \begin{itemize}
        \item \verb|axis equal;|
    \end{itemize}
    \item \textbf{felület rajza}
    \begin{itemize}
        \item \verb|surf(x3, y3, z3)|
    \end{itemize}
    \item \textbf{3d paraméteresen megadott görbe}
    \begin{itemize}
        \item \verb|plot3()|
    \end{itemize}
    \item \textbf{fájlkezelés}:
\begin{minted}[breaklines]{matlab}
f=fopen('d2.txt','r'); % r: olvasás, w: írás, f a file "kezelő", int
disp(f);
adatok=fscanf(f,'%f %f %f\n',[3 inf]); % 3 adat van a sorban, összes sor olvas
xyz=adatok'; % ' a transzponáltja
fclose(f);
x=adatok(1,:)'; % adatok 1. sora x vektorba, transzponálva
y=adatok(2,:)';
z=adatok(3,:)';
\end{minted}
\end{itemize}